\section{Additivity with a finite-dimensional partner}\label{sec:partners}

The branch estimate in \Cref{lem:partner-branch} applies to every
finite-dimensional partner. We now use it to prove the three assertions
of \Cref{cor:partner-additivity}. For the Holevo bound, the key point is
that the conditional partner outputs form an ensemble whose mean is
exactly the partner marginal of the mean full output.

\begin{proof}[Proof of Corollary~\ref{cor:partner-additivity}]
For the Holevo upper bound, it suffices to consider pure-state ensembles.
Every mixed input has a finite spectral decomposition. Refining each
signal into its pure components leaves the mean output unchanged and,
by entropy concavity \eqref{eq:concavity}, does not increase the mean
individual output entropy. Thus refinement cannot decrease Holevo
information; see also Schumacher--Westmoreland
\cite{SchumacherWestmoreland1997}.
The refined pure inputs may still be entangled between the two factors.

Fix a finite ensemble $\{q_a,\proj{\psi_a}\}$ of such inputs.
For each signal use
the decomposition in Lemma~\ref{lem:partner-branch}, with population $p_a$
and branch vectors $\psi_{a,0},\psi_{a,1}$. Put
\[
 \begin{gathered}
 \overline p=\sum_aq_ap_a,\qquad
 w_{a,0}=q_a(1-p_a),\qquad w_{a,1}=q_ap_a,\\
 \tau_{a,i}=\mathcal N(\proj{\psi_{a,i}}).
 \end{gathered}
\]
Here and below $i\in\{0,1\}$.
The nonnegative weights $w_{a,i}$ sum to one. If
$\overline\rho=\sum_aq_a\proj{\psi_a}$ and $\overline\rho_B$ is its
partial trace over the first input qubit, then
\begin{equation}\label{eq:partner-marginal}
 \begin{aligned}
 \overline\rho_B&=\sum_{a,i}w_{a,i}\proj{\psi_{a,i}},\\
 \overline\tau:=\sum_{a,i}w_{a,i}\tau_{a,i}
 &=\mathcal N(\overline\rho_B)
 =\Tr_{\C^2}\bigl((\A\otimes\mathcal N)(\overline\rho)\bigr).
 \end{aligned}
\end{equation}
The cross terms disappear in the first identity because
$\langle0|1\rangle=0$, regardless of any overlap between the branch
vectors. The last identity uses trace preservation of $\A$. Thus the
branch ensemble for $\mathcal N$ has exactly the second marginal of
the mean full output.

Write $\sigma_a=(\A\otimes\mathcal N)(\proj{\psi_a})$ and
$\overline\sigma=\sum_aq_a\sigma_a$. The first output marginal has
excited-state population $t(\overline p)=(1-\gamma)\overline p+\gamma\nu$;
here trace preservation of $\mathcal N$ identifies that marginal with
$\A$ applied to the first input marginal. Apply the diagonal-entropy
bound \eqref{eq:diagonal} in the product of the first output's
computational basis and an eigenbasis of $\overline\tau$. The resulting
joint distribution has marginals $(1-t(\overline p),t(\overline p))$
and the eigenvalue distribution of $\overline\tau$. Classical
subadditivity, as proved in \Cref{sec:holevo}, gives
\begin{equation}\label{eq:partner-mean-upper}
 \ent(\overline\sigma)
 \le h(t(\overline p))+\ent(\overline\tau).
\end{equation}
On the other hand, Lemma~\ref{lem:partner-branch} and the convexity in
Lemma~\ref{lem:e-convex} imply
\begin{equation}\label{eq:partner-ensemble-lower}
 \begin{split}
 \sum_aq_a\ent(\sigma_a)
 &\ge\sum_aq_ae(p_a)+\sum_{a,i}w_{a,i}\ent(\tau_{a,i})\\
 &\ge e(\overline p)+\sum_{a,i}w_{a,i}\ent(\tau_{a,i}).
 \end{split}
\end{equation}
Subtracting \eqref{eq:partner-ensemble-lower} from
\eqref{eq:partner-mean-upper} yields
\begin{align*}
 \ent(\overline\sigma)-\sum_aq_a\ent(\sigma_a)
 &\le h(t(\overline p))-e(\overline p)\\
 &\quad+\ent(\overline\tau)-\sum_{a,i}w_{a,i}\ent(\tau_{a,i})\\
 &\le (\ln2)\bigl(\chi(\A)+\chi(\mathcal N)\bigr).
\end{align*}
For the last line, the one-use case of \Cref{thm:main} bounds the scalar
term, and \eqref{eq:holevo} bounds the Holevo information of the branch
ensemble, whose mean is \eqref{eq:partner-marginal}. Taking the
supremum proves the upper bound in \eqref{eq:partner-holevo}.
For the reverse bound, take the phase-pair ensemble attaining
$\chi(\A)$ and, for each $\varepsilon>0$, a finite ensemble with
Holevo information greater than $\chi(\mathcal N)-\varepsilon$.
Their product ensemble has the sum of their Holevo informations,
because both the mean outputs and individual outputs are tensor
products. Taking the supremum proves \eqref{eq:partner-holevo}.

For minimum output entropy, concavity \eqref{eq:concavity} shows that
some pure input attains the minimum of any finite-dimensional channel:
a pure-state decomposition cannot have every component output entropy
larger than that of the mixed input. For $\A$, the pure qubit calculation
\eqref{eq:single-entropy} therefore gives
\begin{equation}\label{eq:minimum-output-value}
 S_{\min}(\A)=\min_{0\le p\le1}e(p)
              =g\bigl(\gamma\nu(1-\nu)\bigr).
\end{equation}
Indeed, the completed square in \eqref{eq:v-main} is minimized at
$p=\nu$, and \Cref{lem:binary-spectrum} implies that $g$ is
nondecreasing. Any pure qubit state with that population attains the
minimum, including at the parameter endpoints.
The right-hand side of \eqref{eq:partner-branch} is at least
$S_{\min}(\A)+S_{\min}(\mathcal N)$. Hence the same is true of every
output entropy of every pure input for $\A\otimes\mathcal N$, and concavity
extends the lower bound to mixed inputs. Tensoring minimizing input
states for the two factors gives the reverse inequality, since entropy
is additive on product states. This proves \eqref{eq:partner-min}.

Finally, a permutation of tensor factors identifies
$(\A\otimes\mathcal N)^{\otimes n}$ with
$\A^{\otimes n}\otimes\mathcal N^{\otimes n}$ up to input and output
unitaries, which do not change $\chi$. Repeated application of
\eqref{eq:partner-holevo} gives, for every integer $n\ge1$,
\begin{equation}\label{eq:partner-blocks}
 \chi\bigl((\A\otimes\mathcal N)^{\otimes n}\bigr)
 =n\chi(\A)+\chi(\mathcal N^{\otimes n}).
\end{equation}
At each step the partner is a tensor product of the remaining copies
of $\A$ with $\mathcal N^{\otimes n}$, hence is again finite dimensional
and completely positive trace preserving. Applying
\eqref{eq:regularization} to \eqref{eq:partner-blocks} gives
\[
 C(\A\otimes\mathcal N)
 =\chi(\A)+\sup_{n\ge1}\frac1n\chi(\mathcal N^{\otimes n})
 =C(\A)+C(\mathcal N),
\]
where the last equality uses $C(\A)=\chi(\A)$ from \Cref{thm:main}.
This proves \eqref{eq:partner-capacity}. In particular, no restriction
has been placed on entanglement among the inputs to the partner's
copies in \eqref{eq:partner-blocks}.
\end{proof}
